\documentclass[11pt]{article}

\usepackage[margin=1in]{geometry}
\usepackage{amsmath,amssymb,amsthm}
\usepackage{booktabs}
\usepackage{array}
\usepackage[hidelinks]{hyperref}
\usepackage{microtype}

\newcommand{\F}{\mathbf F}
\newcommand{\PGL}{\operatorname{PGL}}
\newcommand{\PGammaL}{\operatorname{P\Gamma L}}
\newcommand{\PRS}{\operatorname{PRS}}
\newcommand{\Gal}{\operatorname{Gal}}
\newcommand{\Deep}{\operatorname{Deep}}
\newcommand{\Stab}{\operatorname{Stab}}
\newcommand{\Osc}{\operatorname{Osc}}
\newcommand{\Span}{\operatorname{span}}
\newcommand{\rank}{\operatorname{rank}}
\newcommand{\Sym}{\operatorname{Sym}}
\newcommand{\NRC}{\operatorname{NRC}}

\title{Reed--Solomon Deep Holes and Reconstruction:\\
Results Established through 23 July 2026}
\author{Tavis Rudd}
\date{23 July 2026}

\begin{document}
\maketitle

\begin{abstract}
This note summarizes the established results of a programme on projective
Reed--Solomon deep holes, determinant invariants, and reconstruction from deepest-syndrome
data.  There are three principal theorem packages: a complete all-field
redundancy-three reconstruction theory; a complete projective-semilinear
classification of the deep holes of $\PRS(q-4)$, of redundancy five; and
all-field/high-field classifications in redundancies six and seven.  The
redundancy-seven theorem is complete for $q\geq37$, while an orbit-reduced
calibration below $37$ remains open.  We state the mathematical results,
indicate the proof methods, and distinguish theorem proofs from bounded
computer-assisted classifications.
\end{abstract}

\section{Coding-theoretic framework}

For a projective Reed--Solomon code of length $q+1$ and redundancy $r$, the
parity-check columns form the rational normal curve
\[
 \nu_{r-1}\bigl(\mathbf P^1(\F_q)\bigr)
 \subset \mathbf P^{r-1}(\F_q).
\]
A projective syndrome is deep precisely when it lies in no span of $r-2$
distinct rational points of this curve.  Equivalently, the corresponding
Hankel-kernel linear system of binary forms has no totally split squarefree
member of degree $r-1$.

For redundancy three, let $A$ be a six-arc in $\mathbf P^2(\F_q)$.  Its
complete projective deepest-syndrome locus is
\[
 U(A)=\mathbf P^2(\F_q)\setminus
 \bigcup_{\{a,b\}\subset A}\overline{ab}.
\]
Thus $u\in U(A)$ exactly when adjoining $u$ to a parity-check matrix for
$A$ gives a one-column MDS extension.

\section{Redundancy three: determinant atlases}

Let
\[
 A=(\mathbf P(h_1),\ldots,\mathbf P(h_6))\subset\mathbf P(V),
 \qquad \dim V=3,
\]
and choose a volume form $\omega\in\det(V)^*$.  For $u\in U(A)$ define
\[
 d_{ij}(u)=\omega(u,h_i,h_j).
\]
The following hold in every characteristic, including supports containing
the point at infinity.

\begin{enumerate}
\item Modulo the exact edge-torus action
\[
 d_{ij}\longmapsto \lambda a_i a_jd_{ij},
\]
balanced four-cycle monomials generate the full invariant quotient.  They
separate finite-field torus orbits without any square-root choice.

\item For rank-two syndromes, and support size at least five, the four-cycle
atlas reconstructs the syndrome.

\item The only structural contraction is the rank-one, or conic, locus.  On
that locus all raw balanced atlases coincide; the unique missing datum is
the radical point of the induced alternating form in
$\mathbf P^1(\F_q)\setminus S$.  No higher balanced edge monomial repairs
this contraction.

\item Projection from $u$ induces
\[
 \beta_u(\bar x,\bar y)=\omega(u,x,y)
\]
on $V/Fu$.  The determinant atlas is exactly the Pl\"ucker presentation of
the labelled projected sextic and hence a point of $M_{0,6}$.  There are
explicit inverses both from the full Pl\"ucker array and from three
normalized cross-ratios.
\end{enumerate}

The proof is algebraic invariant theory together with explicit lattice
generation for the edge torus and direct Pl\"ucker inversion.

\subsection{The first sharp rank-one collision}

For
\[
 S=\{0,1,2,3,4,\infty\}\subset\mathbf P^1(\F_{11}),
\]
the full semilinear support stabilizer is $V_4$, with complement orbits
\[
 \{5,10\},\qquad \{6,7,8,9\}.
\]
Rank-one syndromes whose radicals lie in these two orbits have the same
raw all-one atlas but are sharply separated by the radical stabilizer
orbit.  Their legal-continuation counts are respectively $11$ and $7$,
and their continuation graphs are
\[
 K_5\sqcup C_5\sqcup K_1,\qquad K_5\sqcup2K_1.
\]
Nontrivial radical stabilizer is the cardinality-minimal intrinsic
discriminator.  The finite assertions are covered by an atomic certificate
and a structurally independent replay.

\section{Multi-projection reconstruction and Gale association}

Normalize a labelled parent by
\[
\begin{split}
 h_1&=(1,0,0),\quad h_2=(0,1,0),\quad h_3=(0,0,1),\\
 h_4&=(1,1,1),\quad h_5=(1,a,b),\quad h_6=(1,c,d).
\end{split}
\]
Each projected sextic supplies one exact linear compatibility row in
\[
 z=(a,b,c,d,bc,ad).
\]
Two coherent projections leave a cubic surface of generic dimension two;
three leave a plane cubic of generic dimension one.  Four independent
projections give a line whose product-compatibility cubic contains the
universal collision $h_4=h_5=h_6$.  Removing that factor leaves a separable
quadratic.

Consequently, four abstract coherent projections recover exactly a
two-sheeted parent cover, not a rationally unique parent.  Explicit
examples over $\F_{101}$ and $\F_{256}$, with trivial common diagonal
stabilizer, show that the ambiguity persists in odd and even characteristic
and after diagonal $S_6$-unlabelling.

The deck involution is Gale association.  On the normalized chart put
\[
 L_0=-\det(h_4,h_5,h_6),\qquad
 L_1=\det\bigl(\nu_2(h_1),\ldots,\nu_2(h_6)\bigr).
\]
The kernel cubic factors as
\[
 st(L_0s+L_1t).
\]
Thus $L_1=0$ is exactly the conic locus and is the reduced branch divisor
in every characteristic.  The equation $L_0=0$ is instead an arc-boundary
collinearity.  Every additional abstract coherent projection preserves the
same Gale pair and cannot choose a sheet.  The proof uses exact elimination,
determinant identities, Gale duality, symbolic certification, and an
independent Gale replay.

\subsection{Semilinear descent}

All gauges, compatibility equations, residual cubics, Gale transforms,
branch divisors, and complete-child incidence cuts commute with
coefficientwise Frobenius.  Off the conic divisor, the only geometric
descent obstruction is the $C_2$ Gale-sheet class:
\begin{itemize}
\item a Kummer square class in odd characteristic;
\item an Artin--Schreier trace bit in characteristic two.
\end{itemize}
Connected gauge groups contribute no additional obstruction.  If $H$ is
the common diagonal stabilizer, an unlabelled sheet descends exactly when
the Frobenius sheet bit lies in the image of
\[
 H\longrightarrow C_2.
\]
A surjective image identifies the sheets and sacrifices uniqueness.  Over
$\F_{q^m}$ the Gale bit is multiplied by $m$: odd extensions preserve it,
and even extensions split it.

The exceptional order-three colour collapse over $\F_8$ is separate:
Frobenius acts freely as two $3$-cycles on the six colours.  This is a
lossy $C_3$ quotient, not a Gale obstruction or a failure of Hilbert 90.
The proof combines explicit frame-overlap formulae, finite-field Kummer
and Artin--Schreier theory, Hilbert 90 for the connected frame gauges, and
a finite-stabilizer image criterion.

\section{Complete-child reconstruction}

Retaining the literal complete child $L=U(A)$ adds information absent from
abstract projections.  Parent recovery from selected centres $T\subset L$
is exactly injectivity of the coherent signature map $\Sigma_T$,
equivalently the intersection of the residual cubic family with the
ambient incidence condition
\[
 U(A')=L.
\]

The all-field result is as follows.
\begin{itemize}
\item For $q\geq16$, the complete child alone determines the unlabelled
six-arc.
\item The same child-rigidity holds for every $q=13$ fibre.
\item In every remaining nonempty recoverable fibre, at most three centres
suffice.
\item The only failures are empty children and the $q=7$ two-point
conic-complement child.  In the latter case all $294$ literal parents have
the same coherent signature even after both available centres are used.
\end{itemize}

The complete nontrivial base-size table is
\[
\begin{array}{c@{\quad}c@{\quad}c@{\quad}c}
\toprule
q&|L|&\text{literal parents}&\text{minimum centres}\\
\midrule
8&4&10&3\\
9&6&8&3\\
9&7&2&2\\
9&8&4&2\\
9&8&2&2\\
11&12&22&3\\
\bottomrule
\end{array}
\]
and every other nonempty residual fibre through $q=13$ is a singleton.

For $q\geq16$ the proof is a secant-union rigidity argument.  For small
fields it is an exhaustive classification of child-derived fifteen-line
decompositions followed by an exact collision hypergraph.  For every pair
of parents and every one of the $720$ diagonal label transporters, the
calculation records the set of centres that distinguish them and computes
the exact transversal number.  An independent implementation recomputes
every emitted parent, atlas, transporter comparison, and hitting set.

\subsection{Exceptional non-GRS controls}

There are exactly four semilinear classes of non-GRS six-arcs with
nonempty deepest-syndrome locus contained in a conic, at
\[
 q=8,9,9,11,
\]
with child sizes $4,6,7,12$ and projective deep-hole orbit profiles
\[
 4,\qquad6,\qquad1+6,\qquad12.
\]
A deletion-trace invariant, obtained by intersecting the unique conic
through each five-point deletion with the complete child, recovers the
parent in the six-point $q=9$ and twelve-point $q=11$ cases, but not in
the $q=8$ or seven-point $q=9$ cases.

For the recovering six-point $q=9$ fibre, the eight parents form two
sheets of four perfect matchings of $K_6$ minus a perfect matching.  Their
shared-trace graph is
\[
 K_{4,4}\setminus4K_2\cong Q_3.
\]
The order-$48$ fixed-child group is the full cube automorphism group.  This
example passes decorated recovery but fails the proposed modular
refinement: the relevant Gram rank is $3$, and the Sylow-$C_3$
augmentation module is projective and hence zero in the stable category.
The stronger coherent determinant atlas recovers all four exceptional
families with exact minimum centre counts
\[
 3,\quad3,\quad2,\quad3.
\]

\section{Redundancy five: complete classification of \texorpdfstring{$\PRS(q-4)$}{PRS(q-4)}}

For every prime power $q\geq7$,
\[
 \rho\bigl(\PRS(q-4)\bigr)=4.
\]
This is unconditional, by the Seroussi--Roth completeness range for the
quartic normal rational curve.

For a syndrome $f=(a_0:\cdots:a_4)$, let
\[
 H_f=
 \begin{pmatrix}
 a_0&a_1&a_2&a_3\\
 a_1&a_2&a_3&a_4
 \end{pmatrix}.
\]
Off the curve, $\ker H_f$ is a pencil of binary cubics.  A
characteristic-free recurrence and confluent-Vandermonde argument proves
\[
 f\text{ is deep}
 \quad\Longleftrightarrow\quad
 \ker H_f\text{ has no totally split squarefree cubic}.
\]

The complete $\PGammaL_2$ classification consists of:
\begin{itemize}
\item the tangent family, with $q(q+1)$ points;
\item the conjugate $\sigma$-secant family, with
\[
 \frac{q(q+1)(q-1)}2
\]
points;
\item if $q\equiv2\pmod3$, rational osculating-pair intersections, with
$q(q+1)/2$ points;
\item if $q\equiv1\pmod3$, conjugate osculating-pair intersections, with
$q(q-1)/2$ points;
\item in characteristic $3$, the common nucleus
\[
 (0:0:1:0:0)
\]
of all osculating planes, together with one wild Artin--Schreier orbit of
size $(q^2-1)/2$ and stabilizer $2q$, selected by a nonsquare condition.
\end{itemize}

The arithmetic mechanism is exact.  Tame cyclic degree-three covers are
deep precisely when their order-three deck transformation is
Frobenius-twisted.  In characteristic three, the corresponding condition
is that the additive map $z\mapsto z^3+az$ have irrational kernel.

For geometrically $S_3$ covers, the off-diagonal fibre-square curve has
bidegree $(2,2)$, arithmetic genus one, and is absolutely irreducible.
Aubry--Perret gives
\[
 \#Y(\F_q)\geq q-2\sqrt q.
\]
After deleting at most twelve diagonal and branch points, a split cubic
exists for every $q\geq23$.  Hence no sporadic orbit survives above $19$.

\subsection{Complete sporadic table}

The sporadic classification is organized by the Frobenius type of the
degree-four branch divisor and, in the fully split case, by its
$j$-invariant:
\[
\begin{array}{>{\raggedright\arraybackslash}p{0.27\textwidth}
              >{\centering\arraybackslash}p{0.13\textwidth}
              >{\raggedright\arraybackslash}p{0.43\textwidth}}
\toprule
\text{branch type}&\text{stabilizer}&\text{fields and orbit sizes}\\
\midrule
$1+1+1+1$, $j=0$&
$A_4$&
$7:28,\ 13:182,\ 19:570$\\
$1+1+1+1$, $j\ne0$&
$V_4$&
$9:180,\ 11:330,\ 13:546,\ 17:1224$\\
$1+1+2$&
$C_2$&
$7:2\mathbin{\times}168,\ 8:3\mathbin{\times}252,\
9:2\mathbin{\times}360,\ 11:660$\\
$2+1+1$ with a cusp&
$C_2$&
$7:168$\\
$1+3$&
$C_3$&
$7:112$\\
\bottomrule
\end{array}
\]
The $A_4$ cases are exactly equianharmonic branch tetrads, characterized
by $j=0$, equivalently by vanishing of the binary-quartic invariant $I$.
The three size-$252$ orbits over $\F_8$ form a free
$\Gal(\F_8/\F_2)$-torsor and fuse semilinearly.

These sporadic configurations persist geometrically over larger fields
but cease to be deep because split fibres appear.  Their geometry is
uniform; their deepness is a bounded-field arithmetic phenomenon.

The proof combines characteristic-free Hankel apolarity, confluent
Vandermonde arguments, stratification by pencil gcd and degree-three
monodromy, finite-field covering-curve estimates, and a two-implementation
exhaustive census for every prime power $7\leq q\leq49$.  A
claim-specific literature audit found no pre-emption of the new families
or completeness theorem, subject to its stated database-coverage caveat.

\section{Redundancy six: all-field existence and orbit counts}

For $\PRS(q-5)$, the covering radius is $5$ for every $q\geq7$.  A
syndrome is deep exactly when its Hankel-kernel net of binary quartics has
no totally split squarefree quartic.

The persistent stratum has a quadratic, rather than growing-degree,
common factor and contains exactly
\[
 \frac{q(q+1)^2}{2}
\]
points.  It consists of tangent fibres, whose common quadratic is a
rational square, and $\sigma$-fibres, whose common quadratic is
irreducible.

For tangent fibres, if $p\ne5$ there is one orbit with
\[
 (|O|,|\Stab|)=\bigl(q(q+1),q-1\bigr).
\]
In characteristic five it splits into
\[
 (q+1,q(q-1)),\qquad(q^2-1,q).
\]
For an irreducible quadratic, identify the fibre with the norm-one torus
$T$.  Its stabilizer acts by $z\mapsto z^5$ and inversion.  If
$5\nmid q+1$ there is one orbit.  If $5\mid q+1$, then
$T/T^5\cong C_5$ modulo inversion gives three orbits with stabilizers
$10,5,5$.  Frobenius acts on $C_5$ by multiplication by $p$.

The complete trivial-gcd exceptional inventory is
\[
\begin{array}{c@{\quad}ccccc}
\toprule
q&7&8&9&11&13\\
\midrule
\#\PGL_2\text{-orbits}&18&11&4&2&1\\
\bottomrule
\end{array}
\]
together with one $q+1$ point characteristic-two nucleus orbit for every
\[
 q=2^m,\qquad m\text{ odd},\quad m\geq5,
\]
and no other trivial-gcd orbits.

The nucleus family is represented by
\[
 \langle1,t,t^4\rangle
\]
and is deep exactly when $m$ is odd.  One odd-characteristic echo occurs
at $q=11$, where six collision lines exhaust the twelve rational points
of the relevant conic.  The second $q=11$ orbit has normal form
\[
 W=\langle u,tu,u^2+5\rangle,\qquad u=t^2-2,
\]
with all fifteen rational split-quadratic factor candidates sharing a
root.  The remaining small-field orbit tables are completely certified,
although intrinsic normal forms for every orbit at $q=7,8,9,13$ have not
yet been compressed into a final conceptual list.

The large-field proof uses the first-polar line
\[
 r\longmapsto(a_1-ra_0,\ldots,a_5-ra_4)\subset\mathbf P^4
\]
to reduce quartic-net splitting to the redundancy-five cubic-pencil
classification.  Maximal minors of the quintic catalecticant bound
intersections with the secant cubic by $3$; the tame cyclic locus costs
at most $4$; pointed $g^2_4$ ramification costs at most $6$.  In
characteristics three and two, possible contained polar lines are
classified separately.

The computational component is an elementwise exhaustive census in
eleven fields, with independent definition-based and Hankel-based scans,
complete orbit closure, stabilizer and Frobenius checks, and a separately
implemented replay.

\section{Redundancy seven: interim theorem}

Let $f=(a_0:\ldots:a_6)$ be a binary sextic syndrome and
\[
 H_f=
 \begin{pmatrix}
 a_0&a_1&a_2&a_3&a_4&a_5\\
 a_1&a_2&a_3&a_4&a_5&a_6
 \end{pmatrix}.
\]
For $r\in\mathbf P^1(\F_q)$, contraction gives
\[
 b(r)=(a_1-ra_0,\ldots,a_6-ra_5),
\]
with the evident value at infinity, and
\[
 H_f((t-r)g)=H_{b(r)}g.
\]
The contractions trace a first-polar line in redundancy-six syndrome
space.  Thus $f$ is non-deep exactly when some contraction has a totally
split squarefree quartic avoiding the marked root $r$.  The avoidance
condition is essential.

Let
\[
 C_f=
 \begin{pmatrix}
 a_0&a_1&a_2&a_3&a_4\\
 a_1&a_2&a_3&a_4&a_5\\
 a_2&a_3&a_4&a_5&a_6
 \end{pmatrix}.
\]
The polar line lies entirely in the redundancy-six catalecticant secant
closure if and only if $\rank C_f\leq2$.  After removing rank-one points
and rational split secants, this leaves exactly the persistent tangent
and conjugate-$\sigma$ stratum, again with
\[
 \frac{q(q+1)^2}{2}
\]
points.

\subsection{Persistent orbit law}

The complete orbit law is obtained from $T/T^6$ modulo inversion and
Frobenius.  Put $d=\gcd(6,q+1)$.  An inversion-fixed class gives
\[
 (|O|,|\Stab|)
 =
 \left(\frac{q(q^2-1)}{2d},2d\right),
\]
whereas a paired class gives
\[
 (|O|,|\Stab|)
 =
 \left(\frac{q(q^2-1)}d,d\right).
\]
Frobenius acts by multiplication by $p$ on $C_d$.  On tangent fibres,
if $p\nmid6$ there is one orbit
\[
 (q(q+1),q-1).
\]
In characteristics two and three there are two:
\[
 (q+1,q(q-1)),\qquad(q^2-1,q).
\]
These formulae have been independently replayed at
$q=4,5,7,8,9$, realizing $d=1,2,3,6$ and both modular cases.

\subsection{The central characteristic-two point}

The redundancy-six nucleus line contracts to the single central syndrome
$e_3$.  Its web is
\[
 \langle1,t,t^4,t^5\rangle.
\]
It is deep exactly over $\F_{2^m}$ with $m$ odd.  For even $m$, the
polynomial $t^4+t$, together with the point at infinity, gives a split
quintic on $\mathbf P^1(\F_4)$.

\subsection{High-field theorem}

For every prime power $q\geq37$,
\[
\Deep\bigl(\PRS(q-6)\bigr)
=
\{\text{persistent tangent/$\sigma$ stratum}\}
\ \cup\
\begin{cases}
\{e_3\},&q=2^m,\ m\text{ odd},\\
\varnothing,&\text{otherwise}.
\end{cases}
\]
Consequently the number of deep syndrome directions is
\[
 \frac{q(q+1)^2}{2}
 +
 \begin{cases}
 1,&q=2^m,\ m\text{ odd},\\
 0,&\text{otherwise}.
 \end{cases}
\]

The proof has three quantitative ingredients:
\begin{enumerate}
\item a pointed trivial-gcd upgrade of the redundancy-six splitting
theorem, using the genus-one fibre estimate with one additional forbidden
root, valid for $q\geq37$;
\item a bidegree union bound for gcd-one contractions, proving the
existence of a split lift for $q\geq16$ unless the common root equals the
marked contraction point;
\item identification of that equality locus with ramification of a
separable $g^3_5$, of degree at most $8$.
\end{enumerate}
Together with at most three secant-closure parameters and at most one
characteristic-two nucleus parameter, the total exceptional budget is
\[
 3+8+1=12<q+1.
\]

What remains open is exactly the orbit-reduced calibration for $q<37$ and
conceptual normal forms for any additional exceptional webs.  No complete
all-field redundancy-seven classification is asserted.

\section{Proof architecture and present boundary}

The results use four complementary types of proof.
\begin{enumerate}
\item \textbf{Exact algebraic proofs:} determinant factorization, torus
invariant lattices, Pl\"ucker inversion, catalecticant rank criteria,
Hankel apolarity, confluent Vandermonde arguments, Gale duality, and
first-polar contraction identities.

\item \textbf{Arithmetic-geometric proofs:} classification of degree-three
cover monodromy, Kummer and Artin--Schreier descent, norm-one torus orbit
arithmetic, Aubry--Perret bounds on genus-at-most-one fibre-square curves,
ramification-degree bounds, and bidegree union bounds.

\item \textbf{Finite group and incidence proofs:} stabilizer-image descent
criteria, secant-union rigidity, exact-cover reconstruction, and collision
hypergraph transversal calculations.

\item \textbf{Certified bounded classifications:} deterministic finite-field
censuses, exact orbit closure and stabilizer computations, frozen compact
certificates, and independently implemented replays.  Where a replay does
not independently establish exhaustiveness, that trusted boundary is
stated explicitly.
\end{enumerate}

These results do not prove the general Reed--Solomon deep-hole conjecture.
What is complete is:
\begin{itemize}
\item the all-field redundancy-three reconstruction theory, including its
Gale ambiguity, semilinear descent, and exact small-field exceptions;
\item the all-field projective-semilinear deep-hole classification at
redundancy five;
\item the all-field existence and orbit-count classification at redundancy
six, with some small exceptional orbits still awaiting conceptual normal
forms;
\item the redundancy-seven classification for $q\geq37$, together with
the exact persistent orbit law and central characteristic-two exception.
\end{itemize}

\end{document}
